\documentclass[12pt]{article}
\usepackage[margin=1in]{geometry}
\usepackage{enumitem}
\usepackage{titlesec}
\usepackage{parskip}
\usepackage{graphicx}
\usepackage{xcolor}
\usepackage{hyperref}
\usepackage{fontenc}

\title{\textbf{Clark et al 2017 RG Notes}\\
\large Clark, William Roberts, Matt Golder, and Sona N. Golder. ``The British Academy Brian Barry Prize Essay: An Exit, Voice and Loyalty Model of Politics'' [2017].\\
\textit{British Journal of Political Science} 47(4): 719-748.}
\author{}
\date{}

\begin{document}

\maketitle

Models exist to identify patterns through recognizing commonalities across otherwise distinct phenomena

Article provides a formalization and extension of Hirshman's \textit{Exit, Voice, and Loyalty} (EVL) framework --- the options individuals have when facing a deteriorating environment

Three departures of this model from EVLN (n = neglect):

(1) Considers strategic interaction between actors on both sides, rather than just the dissatisfied actor

(2) Collapses ``neglect'' into ``exit''

(3) Focus on institutional arrangement without considering normative / social-psychological foundation of EVLN

``In effect, the decision to respond to a deleterious change with exit, voice, or loyalty is always a political decision. \textbf{Politics does not just begin when voice is chosen; it begins when voice is considered.''}

Formulation of a sequential game, whereby the government can respond to voice (and the citizen can respond again):

\begin{center}
\includegraphics[width=0.85\textwidth]{figure-1.png}
\end{center}

``It is the citizen's credible exit threat that makes the effective use of voice possible. A dependent government responds positively only because it knows that the citizen will exit if it ignores her use of voice.

See Harmel \& Janda 1994 -- discussion of why party change might occur

\end{document}
